\begin{proposition}[7.1.3]
\label{IV.2.s7.1.3-fr}
Sean $A$ y $A'$ dos anillos locales noetherianos,
$\mathfrak m$ el ideal maximal de $A$ y
$\varphi:A\to A'$ un homomorfismo local. Se supone
que $A'$ es un $A$-módulo plano y que $A'$ es
equidimensional y catenario. Entonces:
\begin{enumerate}
\item[\textup{(i)}] $A$ es equidimensional y catenario.
\item[\textup{(ii)}] Supongamos, además, que $\mathfrak mA'$ es
un ideal de definición de $A'$. Entonces, para todo
ideal $\mathfrak a$ de $A$, $A'/\mathfrak aA'$ es
equidimensional si y solo si $A/\mathfrak aA$ lo es.
En particular, para todo ideal primo $\mathfrak p$
de $A$, $A'/\mathfrak pA'$ es equidimensional.
\end{enumerate}
\end{proposition}

\begin{proof}
Demostremos primero la segunda afirmación de (ii).
Pongamos $X=\operatorname{Spec}(A)$,
$X'=\operatorname{Spec}(A')$,
$Y=V(\mathfrak p)=\operatorname{Spec}(A/\mathfrak p)$
y $Y'=V(\mathfrak pA')=\operatorname{Spec}(A'/\mathfrak pA')$.
Si $f:X'\to X$ es el morfismo correspondiente a
$\varphi$, $Y'$ es también el subpreesquema cerrado
$f^{-1}(Y)$ de $X'$. Como $\mathfrak mA'$ es un ideal
de definición de $A'$, se tiene
\begin{equation}
\dim(X)=\dim(X')
\tag{7.1.3.1}
\label{IV.2.eq.7.1.3.1-fr}
\end{equation}
en virtud de (6.1.3). Del mismo modo,
$\mathfrak mA'/\mathfrak pA'$ es un ideal de definición
de $A'/\mathfrak pA'$ y $Y'$ es plano sobre $Y$ (2.1.4);
por tanto, según (6.1.3),
\begin{equation}
\dim(Y)=\dim(Y').
\tag{7.1.3.2}
\label{IV.2.eq.7.1.3.2-fr}
\end{equation}

Sean $y$ el punto genérico de $Y$, $Y'_i$
($1\leqslant i\leqslant r$) las componentes irreducibles
de $Y'$, y $y'_i$ el punto genérico de $Y'_i$
($1\leqslant i\leqslant r$). Se tiene $f(y'_i)=y$
para todo $i$ (2.3.4). Por otra parte $y'_i$ es un
punto maximal de la fibra $f^{-1}(y)$
(\textbf 0$_{\mathrm
I}$, 2.1.8); por consiguiente,
$\mathscr O_{X',y'_i}\otimes_{\mathscr O_{X,y}}\boldsymbol{k}(y)=
\mathscr O_{X',y'_i}/\mathfrak m_y\mathscr O_{X',y'_i}$
tiene dimensión 0. En otras palabras,
$\mathfrak m_y\mathscr O_{X',y'_i}$ es un ideal de definición
de $\mathscr O_{X',y'_i}$; se puede aplicar de nuevo
(6.1.3), que da
\begin{equation}
\dim(\mathscr O_{X',y'_i})=\dim(\mathscr O_{X,y})
\tag{7.1.3.3}
\label{IV.2.eq.7.1.3.3-fr}
\end{equation}
es decir, por (5.1.2),
\begin{equation}
\operatorname{codim}(Y'_i,X')=\operatorname{codim}(Y,X)
\tag{7.1.3.4}
\label{IV.2.eq.7.1.3.4-fr}
\end{equation}
para todo $i$. Como $A'$ es equidimensional y
catenario, en virtud de (\textbf 0, 14.3.5) se tiene
\begin{equation}
\dim(X')=\dim(Y'_i)+\operatorname{codim}(Y'_i,X')
\tag{7.1.3.5}
\label{IV.2.eq.7.1.3.5-fr}
\end{equation}
\oldpage[IV-2]{184}
para todo $i$. De (7.1.3.1), (7.1.3.2) y
(7.1.3.4), junto con la desigualdad
$\dim(Y'_i)\leqslant\dim(Y')$, se deduce
\begin{equation}
\dim(X)\leqslant\dim(Y)+\operatorname{codim}(Y,X)
\tag{7.1.3.6}
\label{IV.2.eq.7.1.3.6-fr}
\end{equation}
y, por (\textbf 0, 14.2.2.2), ambos miembros de esta
desigualdad son iguales. Además, esta igualdad implica
\begin{equation}
\dim(Y'_i)=\dim(Y')=\dim(Y)
\tag{7.1.3.7}
\label{IV.2.eq.7.1.3.7-fr}
\end{equation}
para todo $i$, lo que demuestra la segunda afirmación
de (ii).

Demostremos ahora la primera afirmación de (ii). Sean
$\mathfrak p_j$ ($1\leqslant j\leqslant n$) los ideales
primos de $A$ minimales entre los que contienen a
$\mathfrak a$, y sean $\mathfrak p'_{jh}$ los ideales
primos de $A'$ minimales entre los que contienen a
$\mathfrak p_jA'$ ($1\leqslant h\leqslant m_j$).
Se sabe entonces, por (2.3.4), que los
$\mathfrak p'_{jh}$ ($1\leqslant j\leqslant n$,
$1\leqslant h\leqslant
m_j$ para cada $j$) son también
los ideales primos de $A'$ minimales entre los que
contienen a $\mathfrak aA'$. Por lo demostrado antes,
para cada $j$, las dimensiones de los anillos
$A'/\mathfrak p'_{jh}$ ($1\leqslant h\leqslant m_j$)
son todas iguales; por tanto, coinciden con
$\dim(A'/\mathfrak p_jA')$, que a su vez es igual a
$\dim(A/\mathfrak p_j)$ por (7.1.3.2). Decir que
$A'/\mathfrak aA'$ es equidimensional significa,
por tanto, que las dimensiones de los $A/\mathfrak p_j$
son todas iguales; es decir, que $A/\mathfrak a$ es
equidimensional.

Demostremos, finalmente, (i). Sea $\mathfrak r'$ un
ideal primo de $A'$ minimal entre los que contienen a
$\mathfrak mA'$, y sea $A''=A'_{\mathfrak r'}$.
Como $A''$ es un $A'$-módulo plano, también es un
$A$-módulo plano. Además, $A''$ es equidimensional y
catenario (\textbf{0}, 16.1.4), y $\mathfrak mA''$ es un ideal
de definición de $A''$. Es posible, pues, limitarse
al caso en que $\mathfrak mA'$ es un ideal de definición
de $A'$. Si $\mathfrak p$ y $\mathfrak q\subset\mathfrak p$
son dos ideales primos de $A$, se puede aplicar (ii)
a $A/\mathfrak q$ y a $A'/\mathfrak qA'$, que es
equidimensional y plano sobre $A/\mathfrak q$.
Si $Z=\operatorname{Spec}(A/\mathfrak q)$ e
$Y=\operatorname{Spec}(A/\mathfrak p)$, se tiene
$\dim(Z)=\dim(Y)+\operatorname{codim}(Y,Z)$.
Además, se puede aplicar (7.1.3.6), que muestra que
$\operatorname{codim}(Y,X)=\dim X\setmin\dim Y$.
Como $X$ es equicodimensional, estas relaciones
muestran que es biequidimensional, en virtud de
(\textbf 0, 14.3.3).
\end{proof}
